% Multi-tenant SaaS, RBAC with scopes, the web attacks covered in the course
% (report plan). Kept short: only what a reader needs for the method.
\section{Background}

\subsection{Multi-tenant applications}
All companies' data is stored in one MySQL database. Every company-owned table carries a company identifier, and four layers enforce separation: an automatic query filter, company-aware foreign keys, a per-request company context taken only from the session, and a cross-company test suite run against every endpoint. Records of another company answer \emph{not found}, so their existence is not revealed.

\subsection{Role-based access control with scopes}
Permissions are granted through roles, and each role assignment has a \emph{scope}: the whole company, a department, a team, or one person. Acting on another person also requires a position above them in the reporting chain. Role assignments will also follow a delegation rule: nobody may grant a permission they do not hold. \pending{present tense once the delegation rule is built (backend issue 56); one sentence on why scopes and the reporting chain make access control harder to get right.}

\subsection{Web attacks covered}
\pending{a short paragraph each, with a citation: SQL injection, cross-site scripting, cross-site request forgery, session hijacking; and why memory-corruption attacks (buffer overflow, return-oriented programming) cannot arise in our own Python and JavaScript code, while the compiled code underneath relies on ASLR, stack canaries, and non-executable memory.}
